\documentclass{article}
\usepackage{amsmath}
\usepackage{amssymb}
\usepackage{geometry}
\geometry{margin=1in}

\begin{document}

\title{PhiBase Arithmetic: A Rational Number System for Exact Computation and Structural Design}
\author{}
\date{\today}

\maketitle

\section*{Introduction}

When working with three dimensional lattice points generated by the Golden Ratio, the use of irrational quantities gets in the way of the underlying symatries.  We now introduce a system that separates out the irrational and results in a simple integer based arithmetic.


This document presents \textbf{PhiBase arithmetic}, an exact computational framework built upon the number system $\mathbb{Q}(\phi)$, where $\phi = \frac{1+\sqrt{5}}{2}$ (the golden ratio). PhiBase arithmetic entirely avoids floating-point operations by representing every number as an exact combination of integers and $\phi$.

Each coordinate $(X, Y, Z)$ in three-dimensional space is represented explicitly as a PhiBase pair of integers:
\[
X = (p_x, x_x),\quad Y = (p_y, x_y),\quad Z = (p_z, x_z),
\]
corresponding exactly to:
\[
X = p_x \cdot \phi + x_x,\quad Y = p_y \cdot \phi + x_y,\quad Z = p_z \cdot \phi + x_z.
\]

\section*{PhiBase Arithmetic from Six-Basis Coordinates}

In the established six-basis structural system, each 3D coordinate $(X, Y, Z)$ is derived from a six-vector of integer coefficients $[a,b,c,d,e,f]$. The explicit transformation to Cartesian coordinates is originally given by:
\begin{align*}
X &= \phi(a + b) + (e - f) \\
Y &= (c - d) + \phi(e + f) \\
Z &= (a - b) + \phi(c + d).
\end{align*}

These coordinates inherently fit within the PhiBase system, as they are rational-linear combinations of integers and $\phi$. This direct relationship leads naturally to the PhiBase representation.

\section*{Explicit PhiBase Conversion}

Converting from a six-vector directly to PhiBase form yields integer pairs as follows:
\[
\begin{aligned}
X &= (a+b)\cdot\phi + (e - f) \quad &\rightarrow\quad (p_x, x_x) &= (a+b, e-f) \\
Y &= (e+f)\cdot\phi + (c - d) \quad &\rightarrow\quad (p_y, x_y) &= (e+f, c-d) \\
Z &= (c+d)\cdot\phi + (a - b) \quad &\rightarrow\quad (p_z, x_z) &= (c+d, a-b).
\end{aligned}
\]

Note that each of $p_x, x_x, p_y, x_y, p_z, x_z$ is explicitly an integer, entirely avoiding floating-point computation.

\section*{Example of PhiBase Conversion}

Consider the six-vector example:
\[
[a,b,c,d,e,f] = [2, -1, 3, 0, -1, 1].
\]
The PhiBase Cartesian coordinates become:
\begin{align*}
p_x &= a+b = 2 + (-1) = 1, & x_x &= e - f = -1 - 1 = -2 \\
p_y &= e+f = -1 + 1 = 0, & x_y &= c - d = 3 - 0 = 3 \\
p_z &= c+d = 3 + 0 = 3, & x_z &= a - b = 2 - (-1) = 3.
\end{align*}

Thus, the exact PhiBase representation is:
\[
X = (1, -2),\quad Y = (0, 3),\quad Z = (3, 3),
\]
which explicitly translates to:
\[
X = 1\cdot\phi - 2,\quad Y = 0\cdot\phi + 3,\quad Z = 3\cdot\phi + 3.
\]

\section*{Benefits of PhiBase Arithmetic}

The advantages of adopting the PhiBase representation are substantial:
\begin{itemize}
\item \textbf{Perfect Accuracy}: Eliminates rounding errors completely.
\item \textbf{Simple Integer Arithmetic}: Direct and exact addition/subtraction of integer pairs.
\item \textbf{Memoizable Operations}: Multiplication and division can be efficiently cached for repetitive computations.
\item \textbf{Clear Debugging and Readability}: Explicit integer pairs facilitate intuitive understanding and error tracing.
\end{itemize}

\section*{Implementation and Next Steps}

To effectively leverage PhiBase arithmetic:
\begin{enumerate}
\item Explicitly adopt PhiBase representation for internal computational code.
\item Store and compute each 3D coordinate explicitly as PhiBase integer pairs.
\item Replace floating-point arithmetic entirely with PhiBase pair operations.
\item Defer floating-point conversion strictly for visualization or final external interfaces.
\end{enumerate}

This strategy maintains exactness and robustness, ensuring that computational precision and structural integrity are consistently preserved.
\vfill
\smallskip
\noindent\small © 2025 James A. Hinds. All rights reserved.
\end{document}

